\documentclass[12pt]{article}
\usepackage[utf8]{inputenc}
\usepackage[T1]{fontenc}
\usepackage[margin=1in]{geometry}
\usepackage{setspace}
\usepackage{xcolor}
\usepackage[colorlinks=true, urlcolor=blue!60!black, linkcolor=blue!60!black]{hyperref}
\setlength{\parindent}{2em}
\setlength{\parskip}{1.25em}
\renewcommand{\baselinestretch}{1.0}
\begin{document}
\begin{center}
{\LARGE \textit{Printz v. United States}, 521 U.S. 898 (1997)\par}
\vspace{0.8em}
{\large Case Brief --- Four-Page Edition\par}
\end{center}
\vspace{1em}

\section{Memory Jogger}

The Court struck down the Brady Act's interim mandate requiring local law enforcement officers to conduct handgun background checks, establishing a categorical anti-commandeering rule: Congress may not conscript state or local executive officers to administer or enforce a federal regulatory program (\href{https://www.law.cornell.edu/supct/html/95-1478.ZO.html}{521 U.S.\ 898}).

\section{Detailed Case Facts}

The Brady Handgun Violence Prevention Act of 1993 directed the Attorney General to build a national instant background check system by November 30, 1998, and set interim procedures until then. A dealer had to obtain a sworn ``Brady Form'' from a handgun buyer, forward it to the ``chief law enforcement officer'' (CLEO) of the buyer's residence, and wait five business days unless the CLEO approved sooner. In states without qualifying alternatives, \S\ 922(s)(2) commanded the CLEO to ``make a reasonable effort to ascertain within 5 business days whether receipt or possession would be in violation of the law.'' Other provisions required CLEOs to destroy records of approved sales and to give rejected buyers written reasons on request.

Two county sheriffs serving as CLEOs, Jay Printz of Ravalli County, Montana, and Richard Mack of Graham County, Arizona, objected to being pressed into federal service and sued separately.

\section{Procedural History}

Both district courts held the CLEO background-check mandate unconstitutional but severable, leaving a voluntary system in place. A divided Ninth Circuit reversed, upholding all the interim provisions. The Supreme Court granted certiorari in both cases, heard argument on December 3, 1996, and reversed on June 27, 1997.

\section{Judicial Votes}

Reversed, \textbf{5--4}.
\begin{itemize}
\item \textbf{Majority:} Justice Scalia, joined by Chief Justice Rehnquist and Justices O'Connor, Kennedy, and Thomas.
\item \textbf{Concurring:} Justice O'Connor; Justice Thomas.
\item \textbf{Dissenting:} Justice Stevens, joined by Justices Souter, Ginsburg, and Breyer; Justice Souter; Justice Breyer, joined by Justice Stevens.
\end{itemize}

\section{Holding}

The Federal Government may neither direct the states to address particular problems nor command state or local officers to administer or enforce a federal regulatory program. The rule is categorical, whether or not policymaking discretion is involved, with no case-by-case balancing. The background-check mandate and the implicit duty to accept Brady Forms are unconstitutional; the record-destruction and written-reasons provisions are simply inoperative for non-participating CLEOs. The Court left open whether the dealer obligations were severable and whether purely ministerial reporting requirements are valid.

\section{Analysis of Opinions}

\textbf{Majority (Scalia).} Conceding that ``there is no constitutional text speaking to this precise question,'' Scalia relied on history, structure, and precedent. \emph{History:} early statutes assigned duties to state \emph{judges}, who are bound by the Supremacy Clause, not to executive officers; when Congress wanted state jails to hold federal prisoners, it merely recommended it. Hamilton's statement in \textit{Federalist} No.\ 27 that state officers would be ``incorporated'' meant only that they could not obstruct federal law. \emph{Structure:} dual sovereignty leaves states ``a residuary and inviolable sovereignty,'' and conscripting state police at no federal cost would blur political accountability; transferring execution to thousands of CLEOs would also undermine the unitary executive under Article II. A law violating state sovereignty is not ``proper'' under the Necessary and Proper Clause. \emph{Precedent:} \textit{New York v.\ United States} controlled; the policymaking/implementation distinction would make state officers ``puppets of a ventriloquist Congress.''

\textbf{O'Connor (concurring).} CLEOs could still volunteer, the provisions were temporary, Congress could use the Spending Clause, and the Court did not decide ministerial reporting duties such as reporting missing children.

\textbf{Thomas (concurring).} The Commerce Clause may not reach intrastate gun sales at all, and the Second Amendment may protect an individual right.

\textbf{Stevens (dissenting).} The Commerce and Necessary and Proper Clauses authorize the mandate, the Tenth Amendment adds no limit, and the Supremacy and Oath Clauses bind state officers. The Framers expected federal use of state officers, and \textit{Garcia} and the Unfunded Mandates Reform Act show political safeguards suffice. The rule could cripple responses to emergencies.

\textbf{Souter and Breyer (dissenting).} Souter found \textit{Federalist} No.\ 27 decisive for enlisting state officers as ``auxiliaries.'' Breyer pointed to Switzerland, Germany, and the European Union, where local implementation preserves local autonomy.

\section{Examples: Future Applications}

\begin{enumerate}
\item \textbf{Same-Side (Invalid): Immigration Detainers.} A federal command that local police hold immigrants 48 hours past release, with no funding condition or opt-in, commandeers state officers.
\item \textbf{Same-Side (Invalid): Barring Repeal of Gambling Bans.} A prohibition on states repealing their sports-gambling bans dictates state lawmaking, as in \textit{Murphy v.\ NCAA} (2018).
\item \textbf{Opposite-Side (Valid): Spending Condition.} Grants conditioned on voluntary background checks leave states a genuine choice.
\item \textbf{Opposite-Side (Valid): Generally Applicable Privacy Law.} Restricting both state DMVs and private brokers from selling driver data regulates the state as a market actor, as in \textit{Reno v.\ Condon} (2000).
\item \textbf{Fence-Sitter: Ministerial Reporting.} A mandate that every sheriff send monthly missing-children statistics to the Justice Department tests the question the Court and O'Connor expressly reserved.
\end{enumerate}

\section{Critique}

\textbf{Originalist and textualist views.} Formalists praise \textit{Printz} for enforcing dual sovereignty against unfunded conscription. Others note the irony that Scalia admitted no text addressed the question and relied on structural inference; Gene Nichol and Saikrishna Prakash argue the historical record was ambiguous or favored federal power to enlist state officers.

\textbf{Progressive and doctrinal views.} Vicki Jackson and Evan Caminker criticize the rigid bright line as historically weak and counterproductive, since it pushes Congress toward larger federal bureaucracies and underestimates voters' ability to assign blame. Yet progressive jurisdictions now invoke the doctrine to shield sanctuary policies from federal immigration mandates.

\section{Key Quotations}

\begin{quote}
``We held in New York that Congress cannot compel the States to enact or enforce a federal regulatory program. Today we hold that Congress cannot circumvent that prohibition by conscripting the State's officers directly. The Federal Government may neither issue directives requiring the States to address particular problems, nor command the States' officers, or those of their political subdivisions, to administer or enforce a federal regulatory program. It matters not whether policymaking is involved, and no case by case weighing of the burdens or benefits is necessary; such commands are fundamentally incompatible with our constitutional system of dual sovereignty.''
\end{quote}
\noindent --- Justice Scalia, Opinion of the Court, 521 U.S.\ at 935.

\begin{quote}
``The power of the Federal Government would be augmented immeasurably if it were able to impress into its service--and at no cost to itself--the police officers of the 50 States.''
\end{quote}
\noindent --- Justice Scalia, Opinion of the Court, 521 U.S.\ at 922.

\begin{quote}
``When a `La[w] . . . for carrying into Execution' the Commerce Clause violates the principle of state sovereignty reflected in the various constitutional provisions we mentioned earlier, supra, at 919, it is not a `La[w] . . . proper for carrying into Execution the Commerce Clause,' and is thus, in the words of The Federalist, `merely [an] ac[t] of usurpation' which `deserve[s] to be treated as such.'''
\end{quote}
\noindent --- Justice Scalia, Opinion of the Court, 521 U.S.\ at 923--924.

\end{document}
